% =============================================================================
% Document compilable avec: pdflatex scRAW_Scientific_Paper.tex
% =============================================================================
\documentclass[10pt,twocolumn,a4paper]{article}

% --- Packages ----------------------------------------------------------------
\usepackage[utf8]{inputenc}
\usepackage[T1]{fontenc}
\usepackage[french]{babel}
\usepackage{mathptmx}
\usepackage{amsmath,amssymb}
\usepackage{geometry}
\geometry{a4paper, margin=1.6cm, columnsep=0.7cm}
\usepackage{hyperref}
\usepackage{enumitem}
\usepackage{abstract}
\usepackage{titlesec}


% --- Hyperref setup ----------------------------------------------------------
\hypersetup{
  colorlinks=true,
  linkcolor=black,
  urlcolor=black,
  citecolor=black,
}

% --- Formatting --------------------------------------------------------------
\titleformat{\section}{\large\bfseries}{\thesection.}{0.5em}{}
\titleformat{\subsection}{\normalsize\bfseries}{\thesubsection.}{0.5em}{}

% --- Title -------------------------------------------------------------------
\title{\vspace{-1.2cm}\LARGE\textbf{scRAW: Single-Cell Rare-population Autoencoder Weighting}\\[4pt]
\large Un cadre pour l'identification de populations cellulaires rares}
\author{\normalsize Fabien BIDET, Victoria BOURGEAIS, Loann Giovannangeli, Patricia THÉBAULT}
\date{}

\begin{document}

\twocolumn[
  \begin{@twocolumnfalse}
    \maketitle
    \begin{abstract}
      \noindent Les données de scRNA-seq sont de grande dimension, bruitées et souvent fortement déséquilibrées. Dans ce contexte, les populations cellulaires rares sont difficiles à isoler car leur signal est dominé par celui des groupes majoritaires. Nous présentons ici la configuration retenue de scRAW, un autoencodeur déterministe conçu pour apprendre un espace latent adapté à ce problème. Le modèle utilisé dans ce travail combine une reconstruction MSE, un entraînement en deux phases, une pondération dynamique par cellule, une régularisation locale par Triplet Loss et une correction de batch par DANN. Le clustering final est réalisé directement dans l'espace latent à l'aide de HDBSCAN. Ce document présente une version courte du cadre méthodologique et des premiers résultats obtenus du modèle sur le jeu de données Human Baron Pancreas.
      \vspace{0.5cm}
    \end{abstract}
  \end{@twocolumnfalse}
]

% =============================================================================
\section{Introduction}
% =============================================================================

L'analyse transcriptomique à l'échelle de la cellule unique permet d'observer finement l'hétérogénéité d'un tissu. Cette richesse s'accompagne de difficultés méthodologiques bien connues : la matrice d'expression est creuse, le bruit technique est important et le nombre de gènes reste très supérieur au nombre de dimensions utiles pour la description de la structure biologique. Un bruit supplémentaire peut aussi apparaître lorsque les données proviennent de plusieurs individus, de plusieurs technologies de séquençage ou de plusieurs temps d'acquisition. Ce phénomène correspond aux effets de batch \cite{tran2020}. Il constitue un enjeu important dans la construction de l'espace latent, car il faut conserver l'information biologique utile sans laisser la représentation être dominée par des variations techniques.

Au-delà de ces sources de variabilité, le clustering non supervisé constitue lui-même une étape délicate. En l'absence d'annotations de référence, la partition obtenue dépend fortement des choix de prétraitement, de sélection de gènes, de réduction de dimension et de la méthode de clustering utilisée. Les recommandations de bonnes pratiques et les études comparatives montrent qu'aucune méthode ne domine dans tous les contextes et que les performances varient sensiblement selon le jeu de données et la structure biologique recherchée \cite{luecken2019,duo2020}. Cette difficulté a conduit au développement de nombreuses méthodes dédiées au clustering single-cell ainsi qu'à des benchmarks comparatifs plus récents, comme scCluBench \cite{xu2025}.

Une difficulté particulière concerne l'identification des populations rares. Ces populations peuvent représenter une faible fraction des cellules observées tout en restant biologiquement importantes, ce qui a motivé plusieurs méthodes dédiées en single-cell \cite{grun2015,jiang2016,wegmann2019}. Dans un schéma d'apprentissage standard, leur contribution à la fonction de perte est souvent trop faible pour orienter l'apprentissage de l'espace latent. Elles peuvent alors être absorbées par des groupes majoritaires ou interprétées comme du bruit. Ce biais se retrouve aussi lors de l'évaluation : des métriques classiques comme le NMI, l'ARI ou le coefficient de silhouette reflètent surtout la qualité globale de séparation et la structure des groupes majoritaires. La reconnaissance des petites populations peut donc être sous-estimée.

scRAW a été conçu pour répondre à cette situation. L'idée générale est de construire un espace latent déterministe, puis d'ajuster l'entraînement pour accorder davantage d'importance aux cellules numériquement rares ou localement isolées. Le modèle utilisé ici repose sur une reconstruction MSE, une pondération dynamique, une régularisation locale et une correction de batch par DANN.


% =============================================================================
\section{Méthodologie}
% =============================================================================

\subsection{Vue d'ensemble}

Le pipeline scRAW comprend quatre étapes principales : la préparation de la matrice d'expression, l'apprentissage d'un espace latent, la mise à jour périodique de poids par cellule, puis le clustering final. Une cellule d'entrée $x_i \in \mathbb{R}^{G}$ est d'abord encodée dans l'espace latent, puis reconstruite par le décodeur :
\[
  z_i = f_{\mathrm{enc}}(x_i),
  \qquad
  \hat{x}_i = f_{\mathrm{dec}}(z_i)
\]
Ici, $x_i$ désigne le profil d'expression de la cellule $i$, $G$ le nombre de gènes considérés, $z_i \in \mathbb{R}^{d_z}$ la représentation latente et $\hat{x}_i$ la reconstruction produite par le décodeur. Les fonctions $f_{\mathrm{enc}}$ et $f_{\mathrm{dec}}$ représentent respectivement l'encodeur et le décodeur.

L'architecture retenue est un autoencodeur déterministe à trois couches cachées dans l'encodeur, suivi d'un décodeur symétrique. Les largeurs des couches sont notées $h_1$, $h_2$ et $h_3$, la dimension latente $d_z$ et le taux de dropout $p_{\mathrm{drop}}$. Chaque bloc caché applique une couche linéaire, puis une normalisation, une activation LeakyReLU et un dropout. Une tête adversariale $g_{\mathrm{batch}}$ est aussi reliée au latent pour prédire la variable de batch $b_i$. Dans le code, cette branche utilise un Gradient Reversal Layer, ce qui laisse la prédiction inchangée vers l'avant mais inverse le gradient remontant vers l'encodeur afin de réduire l'information de batch dans $z_i$.

\subsection{Reconstruction et déroulement de l'entraînement}

Dans la configuration retenue, la reconstruction repose sur une erreur quadratique entre la valeur observée et la valeur reconstruite pour chaque gène. Pour chaque cellule, ces erreurs sont ensuite moyennées sur l'ensemble des gènes afin d'obtenir la perte de reconstruction $\mathcal{L}_{\mathrm{recon},i}$.

Pendant l'entraînement, une partie des gènes est masquée aléatoirement à l'entrée. Le taux de masquage est noté $p_{\mathrm{mask}}$. Le modèle ne cherche donc pas seulement à recopier l'entrée, mais à reconstruire le profil d'expression initial à partir d'une information partiellement cachée. Les gènes masqués et non masqués peuvent contribuer différemment à la perte de reconstruction. Cette importance relative est réglée par le coefficient $\omega$ : plus $\omega$ est élevé, plus la reconstruction des positions masquées est favorisée. En l'absence de masque, on retrouve une moyenne quadratique classique sur les gènes.

Ce principe s'inscrit dans la famille des approches à masquage étudiées récemment pour le single-cell, notamment dans scMAE \cite{fang2024}, et plus largement dans les travaux consacrés à l'apprentissage auto-supervisé en génomique single-cell \cite{richter2025}. Dans la configuration retenue, ce masquage reste actif tout au long de l'entraînement.

L'entraînement est divisé en deux temps. Une première phase de warm-up, de durée $E_{\mathrm{warm}}$, laisse toutes les cellules contribuer de manière uniforme afin de stabiliser le latent. Une seconde phase introduit la pondération dynamique, tandis que la correction de batch et la Triplet Loss peuvent s'activer à des époques propres à chacune. La perte optimisée peut alors s'écrire
\[
  \mathcal{L}_{\mathrm{total}}
  =
  \frac{1}{N}\sum_{i=1}^{N} w_i\,\mathcal{L}_{\mathrm{recon},i}
  +
  \rho(e)\,\lambda_{\mathrm{rare}}\,\mathcal{L}_{\mathrm{triplet}}
  +
  \lambda_{\mathrm{adv}}\,\mathcal{L}_{\mathrm{batch}}
\]
Dans cette équation, $\mathcal{L}_{\mathrm{total}}$ désigne la perte optimisée pendant la phase principale, $w_i$ le poids attribué à la cellule $i$, $\lambda_{\mathrm{rare}}$ le coefficient de la Triplet Loss, $\rho(e)$ le coefficient de ramp-up à l'époque $e$, $\mathcal{L}_{\mathrm{triplet}}$ la perte contrastive locale, $\lambda_{\mathrm{adv}}$ le poids du terme de correction de batch et $\mathcal{L}_{\mathrm{batch}}$ la perte associée au discriminateur de batch.

Pendant le warm-up, les poids restent uniformes avec $w_i=1$, mais la branche DANN peut déjà être active si $e \ge e_{\mathrm{adv}}$. Le terme triplet, lui, s'active à partir d'une époque $e_{\mathrm{trip}}$. La durée totale de l'entraînement est notée $E$, la taille des mini-batches $B$ et le taux d'apprentissage $\eta$. Une recherche d'hyperparamètres avec Optuna \cite{akiba2019} a été utilisée pour retenir ces paramètres.

\subsection{Pondération dynamique et Triplet Loss}

Le poids d'une cellule combine deux informations. La première vient de la taille du pseudo-cluster auquel elle appartient. La seconde repose sur son isolement local dans l'espace latent, mesuré à partir de distances de voisinage. Les pseudo-labels $c_i$ sont recalculés périodiquement à partir du latent. Ils sont obtenus par Leiden pendant l'entraînement, en ajustant sa résolution pour obtenir un nombre de groupes aussi proche que possible du nombre connu de labels du jeu de données. Les deux composantes du poids sont ensuite fusionnées de manière additive, normalisées, bornées entre $w_{\min}$ et $w_{\max}$, puis lissées dans le temps par un coefficient de momentum $\mu$. L'intervalle de mise à jour de ces poids est noté $T_w$.

La Triplet Loss agit en parallèle sur la structure locale du latent. Elle s'appuie sur un ensemble d'ancres sélectionnées à partir des poids, puis compare pour chaque ancre une cellule du même pseudo-groupe et une cellule d'un groupe différent. Dans sa forme simple, ce terme peut s'écrire
\[
  \mathcal{L}_{\mathrm{triplet}}
  =
  \frac{1}{|\mathcal{A}|}
  \sum_{a \in \mathcal{A}}
  \max \left(
    0,\,
    \|z_a - z_{p(a)}\|_2
    -
    \|z_a - z_{n(a)}\|_2
    +
    m
  \right)
\]
où $\mathcal{A}$ désigne l'ensemble des ancres retenues, $z_a$ l'embedding latent de l'ancre $a$, $z_{p(a)}$ celui du positif associé, $z_{n(a)}$ celui du négatif associé et $m$ la marge du terme triplet. Les distances utilisées dans ce terme sont des distances euclidiennes dans l'espace latent. Dans le code, le positif retenu est le plus éloigné du même pseudo-groupe, tandis que le négatif retenu suit une règle semi-hard. L'activation de ce terme commence à l'époque $e_{\mathrm{trip}}$ et progresse ensuite de manière graduelle.

\subsection{Clustering final}

En fin d'entraînement, le clustering est réalisé directement dans l'espace latent
\[
  Z=\{z_1,\dots,z_N\}\subset\mathbb{R}^{d_z}
\]
où $Z$ est l'ensemble des embeddings latents en fin d'entraînement, $z_i$ l'embedding de la cellule $i$ et $d_z$ la dimension de l'espace latent. À partir de cet espace, le clustering final est réalisé avec HDBSCAN \cite{mcinnes2017}. Les hyperparamètres associés à cette étape sont notés $m_{\mathrm{cluster}}$ pour la taille minimale de cluster, $m_{\mathrm{samples}}$ pour la densité minimale locale, $s_{\mathrm{sel}}$ pour la règle de sélection des clusters et $r_{\mathrm{noise}}$ pour la stratégie de réassignation du bruit.

% =============================================================================
\section{Résultats}
% =============================================================================



\subsection{Performance globale}



\subsection{Détection des populations rares}

Les visualisations du latent peuvent notamment être présentées sous forme de projections UMAP \cite{becht2019}, en complément des tableaux de performance par classe.

\subsection{Robustesse et sensibilité}



% =============================================================================
\section{Discussion}
% =============================================================================


\begingroup
\small
\begin{thebibliography}{12}

\bibitem{tran2020}
Hoa Thi Nhu Tran, Kok Siong Ang, Marion Chevrier, Xiaomeng Zhang, Nicole Yee Shin Lee, Michelle Goh, and Jinmiao Chen.
\newblock A benchmark of batch-effect correction methods for single-cell RNA sequencing data.
\newblock \textit{Genome Biology}, 21:12, 2020.

\bibitem{luecken2019}
Malte D. Luecken and Fabian J. Theis.
\newblock Current best practices in single-cell RNA-seq analysis: a tutorial.
\newblock \textit{Molecular Systems Biology}, 15(6):e8746, 2019.

\bibitem{duo2020}
Angelo Duò, Mark D. Robinson, and Charlotte Soneson.
\newblock A systematic performance evaluation of clustering methods for single-cell RNA-seq data.
\newblock \textit{F1000Research}, 7:1141, 2020.

\bibitem{xu2025}
Ping Xu, Zaitian Wang, Zhirui Wang, Pengjiang Li, Jiajia Wang, Ran Zhang, Pengfei Wang, and Yuanchun Zhou.
\newblock scCluBench: Comprehensive Benchmarking of Clustering Algorithms for Single-Cell RNA Sequencing.
\newblock \textit{CoRR}, abs/2512.02471, 2025.

\bibitem{grun2015}
Dominic Grün, Anna Lyubimova, Lennart Kester, Kay Wiebrands, Onur Basak, Nobuo Sasaki, Hans Clevers, and Alexander van Oudenaarden.
\newblock Single-cell messenger RNA sequencing reveals rare intestinal cell types.
\newblock \textit{Nature}, 525(7568):251--255, 2015.

\bibitem{jiang2016}
Lan Jiang, Huidong Chen, Luca Pinello, and Guo-Cheng Yuan.
\newblock GiniClust: detecting rare cell types from single-cell gene expression data with Gini index.
\newblock \textit{Genome Biology}, 17(1):144, 2016.

\bibitem{wegmann2019}
Rebekka Wegmann, Marilisa Neri, Sven Schuierer, Bilada Bilican, Huyen Hartkopf, Florian Nigsch, Felipa Mapa, Annick Waldt, Rachel Cuttat, Max R. Salick, Joe Raymond, Ajamete Kaykas, Guglielmo Roma, and Caroline Gubser Keller.
\newblock CellSIUS provides sensitive and specific detection of rare cell populations from complex single-cell RNA-seq data.
\newblock \textit{Genome Biology}, 20:142, 2019.

\bibitem{fang2024}
Zhaoyu Fang, Ruiqing Zheng, and Min Li.
\newblock scMAE: a masked autoencoder for single-cell RNA-seq clustering.
\newblock \textit{Bioinformatics}, 40(1):btae020, 2024.

\bibitem{richter2025}
Till Richter, Mojtaba Bahrami, Yufan Xia, David S. Fischer, and Fabian J. Theis.
\newblock Delineating the effective use of self-supervised learning in single-cell genomics.
\newblock \textit{Nature Machine Intelligence}, 7:68--78, 2025.

\bibitem{akiba2019}
Takuya Akiba, Shotaro Sano, Toshihiko Yanase, Takeru Ohta, and Masanori Koyama.
\newblock Optuna: A Next-generation Hyperparameter Optimization Framework.
\newblock In \textit{Proceedings of the 25th ACM SIGKDD International Conference on Knowledge Discovery and Data Mining}, pages 2623--2631, 2019.

\bibitem{mcinnes2017}
Leland McInnes, John Healy, and Steve Astels.
\newblock hdbscan: Hierarchical density based clustering.
\newblock \textit{Journal of Open Source Software}, 2(11):205, 2017.

\bibitem{becht2019}
Etienne Becht, Leland McInnes, John Healy, Charles-Antoine Dutertre, Immanuel W. H. Kwok, Lai Guan Ng, Florent Ginhoux, and Evan W. Newell.
\newblock Dimensionality reduction for visualizing single-cell data using UMAP.
\newblock \textit{Nature Biotechnology}, 37:38--44, 2019.

\end{thebibliography}
\endgroup

\end{document}
% =============================================================================
